\section{Per-figure walkthrough}
\label{sec:figure-walkthrough}

The body shows $21$ figures and uses each in a single paragraph of
inference. This section is the slow walk-through: for every
figure, we name the mechanism it captures, the precise feature in
the picture that proves the mechanism is firing, the alternative
explanations we considered and rejected, and the operational
consequence that follows. The order is deliberate: data
diagnostics first, then training dynamics, then frontiers, then
landscape. Reading the figures in this order is a faithful
reconstruction of the report's narrative.

\subsection{Figure F-1: Per-client size and base mortality
(Figure~\ref{fig:client-mortality})}
\textbf{Mechanism.} The natural-clients heterogeneity that drives
every other result. Two histograms side by side: client size in
log-rows, client mortality in $\%$.

\textbf{What to look at.} The mortality bar for Neuro Stepdown is
$2.0\%$, the bar for MICU is $14.8\%$ -- a $7\!\times$ ratio.
Combined with the size ratio ($\sim$$15\!\times$), the
\emph{effective number of deaths} per client spans
$5$ (Neuro Stepdown) to $\sim$$1{,}600$ (MICU), a $300\!\times$
range. This is the formal driver of the small-client failure mode.

\textbf{Alternatives ruled out.} ``The variance is just noise'' --
no, the size ratios are deterministic features of the data. ``The
variance won't matter'' -- no, the worst-client recall column in
every method's row is non-zero only when the small clients are
explicitly defended.

\textbf{Operational read.} Any reported aggregate metric is a
weighted average that under-represents the small ICUs by a factor
of $15\!\times$. Always consult the per-client table.

\subsection{Figure F-2: Confusion matrix per ICU
(Figure~\ref{fig:cm})}
\textbf{Mechanism.} The clinical projection of the federated
prediction onto the per-ICU operating point. A $9$-panel grid of
$2\!\times\!2$ matrices.

\textbf{What to look at.} The Neuro Intermediate panel: TP$=4$,
FN$=6$, FP$=8$, TN$=486$. Six false negatives out of ten deaths is
the exact $0.4$ sensitivity number. The bottom-right of every
panel is dominated by TN -- accuracy is high, but accuracy is
uninformative on rare-class data.

\textbf{Alternatives ruled out.} ``The model just hasn't seen
enough deaths.'' The pattern repeats across all four
\fedprox{}-$\mu$ values; what changes between methods is the
\emph{ratio} of FN to TP. \fedprox{}-$0.1$ flips one Neuro
Intermediate FN into a TP -- exactly one death caught instead of
missed. That is real clinical impact at $\sim$$10$ patients per
year on a unit this size.

\subsection{Figure F-3: \cvar{}-$\alpha$ collapse
(Figure~\ref{fig:cvar-alpha-collapse})}
\textbf{Mechanism.} Discreteness of the tail size at $C\!=\!5$.

\textbf{What to look at.} Three points on top of each other at
$\alpha\!\in\!\{0.75, 0.9, 0.95\}$. The visual proof is that the
markers occlude one another exactly. The remaining two points
($\alpha\!=\!0$, $\alpha\!=\!0.5$) are at distinct positions
because the tail size is distinct ($5$ and $2$).

\textbf{Alternatives ruled out.} ``Maybe \cvar{} converges to the
same place anyway.'' No: the tail \emph{set} is identical because
$\lfloor 5(1-\alpha) \rfloor$ collapses to the same $1$. The
mechanism is integer arithmetic, not training dynamics.

\textbf{Operational read.} Don't deploy \cvar{}-$\alpha$ at small
$C$. If \cvar{} is desired, use $C \geq 20$.

\subsection{Figure F-4: Dirichlet trends
(Figure~\ref{fig:dirichlet})}
\textbf{Mechanism.} Heterogeneity stress test.

\textbf{What to look at.} Mean accuracy degrades smoothly with
$\beta\!\to\!0$, but the seed-std band widens dramatically at
$\beta\!=\!0.1$. The \fedprox{} curve is consistently above the
\fedavg{} curve, with the gap widening at lower $\beta$.

\textbf{Alternatives ruled out.} ``It's just convergence.'' No:
even at $20$ rounds the per-seed trace is already qualitatively in
the wide-band regime; the mean settles, the variance does not.

\textbf{Operational read.} On heterogeneous deployments, demand
seed-std as part of the SLA, not just mean.

\subsection{Figure F-5: Drift norms over rounds
(Figure~\ref{fig:drift-norms})}
\textbf{Mechanism.} \fedprox{}'s proximal contraction.

\textbf{What to look at.} The four \fedprox{}-$\mu$ traces are
ordered top-to-bottom by $\mu$ value: $\mu\!=\!0$ at $\sim$$0.015$,
$\mu\!=\!0.001$ at $\sim$$0.011$, $\mu\!=\!0.01$ at $\sim$$0.005$,
$\mu\!=\!0.1$ at $\sim$$0.002$. The gap is largest at the small
clients (per-client decomposition in
\texttt{drift\_norms\_per\_client.csv}).

\textbf{Alternatives ruled out.} ``Drift could just be sampling
noise.'' No: the per-seed std on the drift norm is small relative
to the mean ($\sim$$10\%$), and the rank order across $\mu$ values
is preserved across all $10$ seeds.

\textbf{Operational read.} Drift is the precise signal that the
proximal anchor produces the predicted contraction. Drift norm
should be logged in deployment as a runtime healthcheck.

\subsection{Figure F-6: Fairness frontier
(Figure~\ref{fig:fairness-frontier})}
\textbf{Mechanism.} Worst-recall vs.\ accuracy.

\textbf{What to look at.} \fedprox{}-$0.1$ is the most
north-east point: highest worst-recall \emph{and} highest
\auprc{}. \fedavg{} is south-west. \cvar{} variants cluster in the
middle-south. The frontier is monotone in $\mu$, not in $\alpha$.

\textbf{Alternatives ruled out.} ``\fedprox{} just trades accuracy
for fairness.'' No -- it does not: \fedprox{}-$0.1$'s accuracy is
within $4\%$ of \fedavg{}'s, which is well within the seed-std.

\textbf{Operational read.} The fairness gain is not a trade-off
in the regime we tested. Deploy \fedprox{}-$0.1$.

\subsection{Figure F-7: \fedprox{} $\mu$ sweep
(Figure~\ref{fig:fairness-frontier})}
\textbf{Mechanism.} Monotone improvement of \auprc{} with $\mu$.

\textbf{What to look at.} A four-point line ascending from
$0.636$ to $0.665$. The slope decreases between $\mu\!=\!0.01$ and
$\mu\!=\!0.1$, hinting at a plateau.

\textbf{Alternatives ruled out.} ``Maybe higher $\mu$ would do
even better.'' Possible but not tested; the plateau hint suggests
it would not.

\textbf{Operational read.} $\mu\!=\!0.1$ is the pragmatic
recommendation; $\mu\!=\!0.3$ is the next experiment.

\subsection{Figure F-8: GA convergence
(Figure~\ref{fig:ga-convergence})}
\textbf{Mechanism.} Stochastic search converges by iteration $9$.

\textbf{What to look at.} Two panels: the search history scatter
(left) shows clusters of evaluations at $E\!=\!1, 2, 3$; the
best-so-far trace (right) drops sharply for the first $9$
iterations and plateaus. The plateau value is the operating-point
fitness.

\textbf{Alternatives ruled out.} ``DE just got lucky early.''
Re-runs would be the cleanest test; we did not run them. The
plateau itself is evidence that DE has saturated the search space.

\textbf{Operational read.} $90$ evaluations is enough; $200$+ is
overkill on this search box.

\subsection{Figure F-9: Local-vs-federated per client
(Figure~\ref{fig:local-vs-fed})}
\textbf{Mechanism.} Knowledge transfer from federation.

\textbf{What to look at.} Eight clients show federation outperforming
local-only (left bar lower than right bar); one client (Neuro
Stepdown) shows the opposite. The flip is not because federation
hurts -- it is because local-only on Neuro Stepdown predicts
``no death'' on every test row (recall $0$, accuracy $\sim$$0.98$).

\textbf{Alternatives ruled out.} ``Federation helps everyone.'' No:
on the smallest, lowest-mortality client, federation can't recover
because the rare class is too sparse. Federation helps the
\emph{recall}, not the \emph{accuracy}.

\textbf{Operational read.} Federation's headline gain is on the
clients that have enough data \emph{and} enough rare-class events.

\subsection{Figure F-10: Loss landscape 1D
(Figure~\ref{fig:landscape-1d})}
\textbf{Mechanism.} Validation loss as a function of position
along the init-to-final axis.

\textbf{What to look at.} Smooth bowl in both backbones. The MLP
minimum is at $\alpha\!=\!0.85$, the logistic minimum at
$\alpha\!=\!1.0$. The MLP overshoots; the logistic does not.

\textbf{Alternatives ruled out.} ``The minimum is exactly the
final round.'' No, only for logistic; for MLP it is not.

\textbf{Operational read.} For MLP, terminate $\sim$$15\%$ before
the final round.

\subsection{Figure F-11: Loss landscape 2D
(Figure~\ref{fig:landscape-2d})}
\textbf{Mechanism.} Two-direction landscape.

\textbf{What to look at.} Logistic surface is a bowl; MLP surface
shows a slight ridge in the second direction. The ridge implies
non-axis-aligned curvature.

\textbf{Alternatives ruled out.} ``The ridge is an artefact of
the random direction.'' Possible -- a Hessian-eigenvector direction
would produce a different picture. The directional independence is
not formally tested.

\textbf{Operational read.} The MLP landscape is non-isotropic; do
not rely on grid search alone for $\eta$ tuning.

\subsection{Figure F-12: LP shadow price
(Figure~\ref{fig:lp-shadow})}
\textbf{Mechanism.} Communication shadow price as a function of
budget.

\textbf{What to look at.} A monotone descending curve in
log-log; a sharp drop where the budget exits the active region;
$\lambda$ then drifts toward $0$ exponentially.

\textbf{Alternatives ruled out.} ``Maybe more budget would help in
some regime.'' The LP says no: $\lambda$ stays below $10^{-9}$ once
past the active edge, which is operationally negligible.

\textbf{Operational read.} The deployment budget should be $\sim$
$300$--$430$ Mb (the active edge), not $1$ Gb.

\subsection{Figure F-13: Method-grouped bars
(Figure~\ref{fig:headline-bars})}
\textbf{Mechanism.} Side-by-side comparison of four metrics.

\textbf{What to look at.} The bars do not all peak on the same
method. Centralized peaks accuracy, \fedprox{}-$0.1$ peaks the
other three. The visual is the strongest argument that ``one
metric'' is the wrong way to read the table.

\textbf{Alternatives ruled out.} ``Pick the best on accuracy.''
That gives centralized, which is non-deployable.

\textbf{Operational read.} Use \fedprox{}-$0.1$ for the
\auprc{}/worst-recall axis. Tune the operating threshold per ICU.

\subsection{Figure F-14: Pareto: \auprc{} vs.\ communication
(Figure~\ref{fig:pareto})}
\textbf{Mechanism.} Loss-vs-bytes Pareto frontier.

\textbf{What to look at.} Two clusters: dense MLP at
$\sim$$300$--$700$ Mb, $\sim$$0.63$--$0.67$ \auprc{};
logistic at $\sim$$6$ Mb, $\sim$$0.60$--$0.61$ \auprc{}. The
empty middle is the missing operating regime.

\textbf{Alternatives ruled out.} ``Top-$k$ would fill the
middle.'' Numerically, no: the natural sparsity is too low for
top-$k$ to be a free dividend (Section~\ref{sec:dossier-sparsity}).

\textbf{Operational read.} Pick by bandwidth: edge ICU $\to$
logistic; tertiary centre $\to$ MLP.

\subsection{Figure F-15: Per-client clinical metrics
(Figure~\ref{fig:per-client-clinical})}
\textbf{Mechanism.} Sensitivity, specificity, \auroc{}, \auprc{}
per ICU at fixed threshold.

\textbf{What to look at.} The Neuro Intermediate column shows
sensitivity $0.4$, specificity $0.984$ -- the model's bias
toward ``no death'' is most extreme at this small client.

\textbf{Alternatives ruled out.} ``Tune the threshold per
client.'' Yes, but the deployment must commit to a single
threshold for governance reasons.

\textbf{Operational read.} \fedprox{}-$0.1$ improves Neuro
Intermediate sensitivity to $\sim$$0.55$ -- catching $1$ extra death
out of $10$.

\subsection{Figure F-16: Per-seed \auprc{}
(Figure~\ref{fig:per-seed-auprc})}
\textbf{Mechanism.} Across-seed reliability.

\textbf{What to look at.} \fedprox{}-$0.1$ box is the smallest;
Dirichlet-$0.1$ box is the largest. The visual ratio of widths is
$\sim$$30\!\times$.

\textbf{Alternatives ruled out.} ``\fedprox{} got lucky on these
seeds.'' Across $10$ prime seeds the std is $0.005$; that is the
narrowest band of any method.

\textbf{Operational read.} \fedprox{}-$0.1$ is the most
\emph{reliable} method, not just the highest mean.

\subsection{Figure F-17: Resource time series
(Figure~\ref{fig:resources})}
\textbf{Mechanism.} Memory and CPU profile during the proposal-
alignment run.

\textbf{What to look at.} Memory peaks at $9.28$~GB during the
sparsity audit; CPU pegs $\sim$$80\%$ during \fedprox{}-$0.01$
client iteration; the watchdog warning line is at $34$~GB and is
never crossed.

\textbf{Alternatives ruled out.} ``Memory headroom is fine
because we never hit it.'' The warning line gives $24$~GB of
buffer; on a tighter machine ($24$~GB total) the watchdog would
trigger, which is by design.

\textbf{Operational read.} The deployment budget for this run on
a $48$~GB box is $\sim$$10$~GB; on a $32$~GB box, the run still
fits but with the watchdog active.

\subsection{Figure F-18: ROC and PR curves
(Figure~\ref{fig:roc-pr})}
\textbf{Mechanism.} Threshold sweep.

\textbf{What to look at.} ROC is Hochrechnung-classic; PR is
the operationally informative one. The PR curve shows
\fedprox{}-$0.1$'s gain is at \emph{all} threshold values, not
just at $0.5$.

\textbf{Alternatives ruled out.} ``The gain is just at the
default threshold.'' No: the PR curve dominates \fedavg{}'s
across the entire threshold range.

\textbf{Operational read.} Deployment threshold tuning will move
the operating point, but not the recommendation.

\subsection{Figure F-19: Search history scatter
(Figure~\ref{fig:ga-scatter})}
\textbf{Mechanism.} GA exploration in $(E, C, \eta)$.

\textbf{What to look at.} The $90$ evaluations fan out in the box;
the best fitness corner is the lower-left of the loss-vs-comm
plane.

\textbf{Alternatives ruled out.} ``GA didn't explore enough.'' The
scatter shows wide coverage; the plateau in the convergence trace
shows convergence.

\textbf{Operational read.} The optimal $(E, C, \eta)$ point is at
$(1, 5, 0.0061)$; deploy this.

\subsection{Figure F-20: Size vs.\ mortality scatter
(Figure~\ref{fig:size-vs-mortality})}
\textbf{Mechanism.} The two-axis client characterisation.

\textbf{What to look at.} Neuro Stepdown is the small-low-mortality
extreme corner; MICU is the large-high-mortality opposite corner.

\textbf{Alternatives ruled out.} ``All clients are similar
enough.'' The scatter spans two orders of magnitude in size and
$7\!\times$ in mortality.

\textbf{Operational read.} The natural-clients axis is the
strongest possible test of fairness; synthetic Dirichlet partitions
are stress tests, not replacements.

\subsection{Figure F-21: Sparsity compression
(Figure~\ref{fig:sparsity})}
\textbf{Mechanism.} Top-$k$ compression ratios across $k$ values.

\textbf{What to look at.} Compression scales linearly with $k$:
$1\%$ gives $\sim$$50\!\times$ savings, $10\%$ gives
$\sim$$5\!\times$. The natural-sparsity bar at the right is at
$\sim$$96\%$ -- almost all entries are non-zero.

\textbf{Alternatives ruled out.} ``Maybe later in training
sparsity would emerge.'' The CSV's per-round breakdown shows
sparsity stays in the $95.8$--$97.7\%$ band throughout; the
gradient does not concentrate.

\textbf{Operational read.} Use compression only when bandwidth is
the binding constraint; on \mimic{}'s budgets it is not.

\paragraph{Closing remark.}
Twenty-one figures, twenty-one one-paragraph mechanism stories.
Each picture maps to a single CSV (cited in the figure docstrings
in \texttt{build\_plots.py}); each picture proves a single claim
in the body. A reader who has walked the body and this section
has seen the entire chain from data to recommendation.
